
\documentclass{article}
\usepackage{colortbl}
\usepackage{makecell}
\usepackage{multirow}
\usepackage{supertabular}

\begin{document}

\newcounter{utterance}

\centering \large Interaction Transcript for game `cladder', experiment `full\_v1.5\_default', episode 9295 with reflexion{-}Qwen3.5{-}27B.
\vspace{24pt}

{ \footnotesize  \setcounter{utterance}{1}
\setlength{\tabcolsep}{0pt}
\begin{supertabular}{c@{$\;$}|p{.15\linewidth}@{}p{.15\linewidth}p{.15\linewidth}p{.15\linewidth}p{.15\linewidth}p{.15\linewidth}}
    \# & \multicolumn{2}{c}{Player} && \multicolumn{2}{c}{Game Master} \\
    \hline

    \theutterance \stepcounter{utterance}  
    & & & \multicolumn{4}{p{0.6\linewidth}}{
        \cellcolor[rgb]{0.9,0.9,0.9}{
            \makecell[{{p{\linewidth}}}]{
                \texttt{\tiny{[P1$\langle$GM]}}
                \texttt{You have attempted this same task before and did not succeed. These are the plans you wrote after those attempts.} \\
\texttt{Trial \#1: New plan: My previous attempt failed because I answered "No" to the question "Is the chance of glimx smaller when observing zuph?", despite my own calculations proving that \$P(\textbackslash{}text\{glimx\} \textbackslash{}mid \textbackslash{}text\{zuph\}) = 0.3125\$ is indeed smaller than \$P(\textbackslash{}text\{glimx\} \textbackslash{}mid \textbackslash{}neg\textbackslash{}text\{zuph\}) = 0.75\$. The correct answer is "Yes". I mistakenly focused on the causal structure (which suggests a positive influence) and the phenomenon of collider bias, leading me to incorrectly conclude the answer was negative, or I simply misread the final comparison. Next time, I will strictly follow the numerical result of the conditional probability comparison to determine the "Yes" or "No" answer, ensuring the initial word matches the mathematical inequality derived (\$0.3125 < 0.75\$ implies "Yes"). I will also ensure the reasoning explicitly states that the observed statistical association (negative) overrides the intuitive causal path in this specific observational query.} \\
\\ 
\\ 
\texttt{You are an expert in causal inference. The following question is not a typical commonsense query, but rather a meticulously designed question created by a professor specializing in causal inference, intended to assess the students' mastery of the course content.} \\
\texttt{Imagine a self{-}contained, hypothetical world with only the following conditions, and without any unmentioned factors or causal relationships: Zory has a direct effect on zuph and glimx. Zuph has a direct effect on jyka. Jyka has a direct effect on glimx. Zory is unobserved. The overall probability of zuph is 96\%. The probability of not zuph and glimx is 3\%. The probability of zuph and glimx is 30\%. Is the chance of glimx smaller when observing zuph?} \\
\texttt{Start your answer with 'Yes' or 'No', followed by additional reasoning or evidence to support your explanation.} \\
            }
        }
    }
     \\ \\

    \theutterance \stepcounter{utterance}  
    & \multicolumn{4}{p{0.6\linewidth}}{
        \cellcolor[rgb]{0.561,0.737,0.561}{
            \makecell[{{p{\linewidth}}}]{
                \texttt{\tiny{[P1$\rangle$GM]}}
                \texttt{Yes} \\
\\ 
\texttt{To determine if the chance of **glimx** is smaller when observing **zuph**, we must compare the conditional probability \$P(\textbackslash{}text\{glimx\} \textbackslash{}mid \textbackslash{}text\{zuph\})\$ with the conditional probability \$P(\textbackslash{}text\{glimx\} \textbackslash{}mid \textbackslash{}neg\textbackslash{}text\{zuph\})\$.} \\
\\ 
\texttt{**1. Calculate \$P(\textbackslash{}text\{glimx\} \textbackslash{}mid \textbackslash{}text\{zuph\})\$:**} \\
\texttt{We are given:} \\
\texttt{*   \$P(\textbackslash{}text\{zuph\}) = 0.96\$} \\
\texttt{*   \$P(\textbackslash{}text\{zuph\} \textbackslash{}cap \textbackslash{}text\{glimx\}) = 0.30\$} \\
\\ 
\texttt{Using the definition of conditional probability:} \\
\texttt{\$\$P(\textbackslash{}text\{glimx\} \textbackslash{}mid \textbackslash{}text\{zuph\}) = \textbackslash{}frac\{P(\textbackslash{}text\{zuph\} \textbackslash{}cap \textbackslash{}text\{glimx\})\}\{P(\textbackslash{}text\{zuph\})\} = \textbackslash{}frac\{0.30\}\{0.96\} = 0.3125\$\$} \\
\\ 
\texttt{**2. Calculate \$P(\textbackslash{}text\{glimx\} \textbackslash{}mid \textbackslash{}neg\textbackslash{}text\{zuph\})\$:**} \\
\texttt{We are given:} \\
\texttt{*   \$P(\textbackslash{}neg\textbackslash{}text\{zuph\} \textbackslash{}cap \textbackslash{}text\{glimx\}) = 0.03\$} \\
\texttt{*   \$P(\textbackslash{}neg\textbackslash{}text\{zuph\}) = 1 {-} P(\textbackslash{}text\{zuph\}) = 1 {-} 0.96 = 0.04\$} \\
\\ 
\texttt{Using the definition of conditional probability:} \\
\texttt{\$\$P(\textbackslash{}text\{glimx\} \textbackslash{}mid \textbackslash{}neg\textbackslash{}text\{zuph\}) = \textbackslash{}frac\{P(\textbackslash{}neg\textbackslash{}text\{zuph\} \textbackslash{}cap \textbackslash{}text\{glimx\})\}\{P(\textbackslash{}neg\textbackslash{}text\{zuph\})\} = \textbackslash{}frac\{0.03\}\{0.04\} = 0.75\$\$} \\
\\ 
\texttt{**3. Comparison:**} \\
\texttt{Comparing the two results:} \\
\texttt{\$\$0.3125 < 0.75\$\$} \\
\texttt{\$\$P(\textbackslash{}text\{glimx\} \textbackslash{}mid \textbackslash{}text\{zuph\}) < P(\textbackslash{}text\{glimx\} \textbackslash{}mid \textbackslash{}neg\textbackslash{}text\{zuph\})\$\$} \\
\\ 
\texttt{**Conclusion:**} \\
\texttt{The probability of **glimx** occurring is significantly lower (31.25\%) when **zuph** is observed compared to when **zuph** is not observed (75\%). This negative association arises due to the specific structure of the causal graph (a "collider" or "fork" structure involving the unobserved variable **Zory** and the path through **jyka**), which creates a statistical dependency that overrides the intuitive expectation of a positive causal link. Therefore, the chance of glimx} \\
            }
        }
    }
    & & \\ \\

    \theutterance \stepcounter{utterance}  
    & & & \multicolumn{2}{p{0.3\linewidth}}{
        \cellcolor[rgb]{0.9,0.9,0.9}{
            \makecell[{{p{\linewidth}}}]{
                \texttt{\tiny{[GM$|$GM]}}
                \texttt{Yes} \\
            }
        }
    }
    & & \\ \\

    \theutterance \stepcounter{utterance}  
    & & & \multicolumn{2}{p{0.3\linewidth}}{
        \cellcolor[rgb]{0.9,0.9,0.9}{
            \makecell[{{p{\linewidth}}}]{
                \texttt{\tiny{[GM$|$GM]}}
                \texttt{yes} \\
            }
        }
    }
    & & \\ \\

    \theutterance \stepcounter{utterance}  
    & & & \multicolumn{2}{p{0.3\linewidth}}{
        \cellcolor[rgb]{0.9,0.9,0.9}{
            \makecell[{{p{\linewidth}}}]{
                \texttt{\tiny{[GM$|$GM]}}
                \texttt{game\_result = WIN} \\
            }
        }
    }
    & & \\ \\

\end{supertabular}
}

\end{document}
